This survey (re)introduces reinforcement learning to economists. I review the theoretical connections between dynamic programming and reinforcement learning, demonstrating how value iteration, Q-learning, and policy gradient methods are common solution methods to the same class of optimization problems. I then examine applications across several domains, including control problems including pricing and inventory management; structural economic models with high-dimensional state spaces; strategic games in which multi-agent algorithms compute equilibria under imperfect information; bandit problems in which economic structure yields tighter regret bounds; and preference learning. The exposition combines formal theory, practical applications and computational illustration.

Both dynamic programming and reinforcement learning solve the Bellman equation; they differ in the information requirements and the way in which the solution is refined. First, dynamic programming requires knowledge of the transition in the environment and the reward function which allows the reduction of the average Bellman error, reinforcement learning estimates value functions only from sampled transitions (obsserved sets of state, action, reward, next-state) which only allows reduction of the sampled Bellman error \textit{at that state-action pair}. This allows us to improve policies in domains where it is easier to build a simulator than specify the model of the environment and rewards e.g. board games, physics simulators for robots. Second, dynamic programming makes a ``breadth-first'' (accross all states and actions) update of the solution at each sweep, while reinforcement learning makes a ``incremental'' (only for the current state and action) update; this greatly reduces the computational burden and enhances scalability. 

Reduction of average Bellman errors gives dynamic programming a geometric rate of convergence to the optimal solution, while the incremental updates and reduction of sampled Bellman errors, when combined with ``sufficient exploration'' of the state-action space, gives reinforcement learning only sublinear convergence guarantees. This is however, quite sufficient in practise, the scalability attained by only sampling transitions and making incremental updates more than makes up for the slower rates of convergence (and brittleness). 

These approximation methods sacrifice theoretical guarantees. RL algorithms lack the convergence assurances of exact dynamic programming. They exhibit sensitivity to hyperparameters and initialization. They can converge to suboptimal policies without diagnostic indication. This survey presents reinforcement learning as a computationally flexible framework while acknowledging its methodological limitations.

This survey focuses on less-surveyed intersections between reinforcement learning and economics, including the shared theoretical foundations, structural estimation, strategic interaction, bandit problems with economic structure, preference learning, and causal inference. Algorithmic collusion, in which independent pricing algorithms learn to sustain supra-competitive prices \citep{Calvano2020}, is treated in a companion thesis chapter \citep{Rawat2026collusion} and omitted here. Reinforcement learning and deep learning methods for solving macroeconomic models with heterogeneous agents constitute a growing literature with dedicated methodological treatments \citep{AtashbarShi2022}, \citep{MaliarMaliarWinant2021}, and \citep{FernandezVillaverdeNunoPerla2024}. Portfolio optimization, optimal execution, and asset pricing via reinforcement learning form a large body of work surveyed comprehensively elsewhere \citep{HamblyXuYang2023}. The inverse problem of inferring preferences from observed behavior using inverse reinforcement learning and structural estimation is treated in a companion survey \citep{RustRawat2026}.

This survey adopts a unified notation for concepts that developed independently across these traditions. A Markov decision process consists of a state space $\mathcal{S}$, an action space $\mathcal{A}$, a transition kernel $P(s'|s,a)$ giving the probability of reaching state $s'$ from state $s$ under action $a$, a reward function $r(s,a)$, and a discount factor $\gamma \in [0,1)$. A policy $\pi(a|s)$ specifies the probability of selecting action $a$ in state $s$; a deterministic policy is written $\pi(s)$. The state-value function $V^\pi(s) = \mathbb{E}_\pi[\sum_{t=0}^\infty \gamma^t r(s_t, a_t) \mid s_0 = s]$ measures expected discounted return from state $s$ under policy $\pi$; the action-value function $Q^\pi(s,a)$ conditions additionally on the first action. The optimal value functions $V^*$ and $Q^*$ correspond to the policy that maximizes expected return. Value functions are often approximated by parameterized functions: $V(s; \mathbf{w})$ for linear approximations with weight vector $\mathbf{w}$, and $V(s; \theta)$ for general parameterizations including neural networks. Updates proceed by a learning rate $\alpha > 0$. The temporal difference error $\delta_t = r_{t+1} + \gamma V(s_{t+1}) - V(s_t)$ measures the discrepancy between successive value estimates and drives most learning algorithms in this survey.

The survey addresses the forward problem, that is, computing optimal policies given a known or simulated environment. Chapter 1 traces the parallel historical development of dynamic programming and reinforcement learning. Chapter 2 develops unified theory connecting planning and learning. Chapters 3 and 4 apply reinforcement learning to control problems and economic models. Chapter 5 examines strategic games. Chapter 6 addresses bandit problems. Chapter 7 discusses reinforcement learning from human feedback. Chapter 8 connects reinforcement learning to causal inference. Chapter 9 concludes.
